\section{优化基础回顾}

本节课承接上一讲的内容，首先回顾梯度下降优化的基本思想，然后深入讨论 Word2Vec 的更多细节，介绍基于计数的词向量方法（共现矩阵与 GloVe），探讨词向量的评估方式和词义消歧问题，最后引入神经网络分类器的基本概念。

\begin{figure}[H]
\centering
\includegraphics[width=0.95\textwidth]{slides-images/slide-001.jpg}
\caption{本讲大纲：涵盖优化基础、Word2Vec 回顾与扩展、共现计数方法、词向量评估、词义问题、分类器与神经网络入门\protect\footnotemark}
\end{figure}
\footnotetext{Slides 第2页。}

\subsection{梯度下降（Gradient Descent）}

梯度下降是训练词向量（以及所有神经网络）的核心优化算法。其基本思想是：

\begin{enumerate}
    \item 我们有一个损失函数 $J(\theta)$，希望找到使其最小化的参数 $\theta$
    \item 计算 $J(\theta)$ 关于 $\theta$ 的梯度 $\nabla_\theta J(\theta)$，得到``下坡''方向
    \item 沿梯度的反方向走一小步，更新参数
    \item 重复上述过程
\end{enumerate}

\begin{figure}[H]
\centering
\includegraphics[width=0.95\textwidth]{slides-images/slide-003.jpg}
\caption{梯度下降示意：从随机初始值出发，沿梯度反方向逐步走向最小值\protect\footnotemark}
\end{figure}
\footnotetext{Slides 第4页。}

参数更新公式为：

$$\theta^{\text{new}} = \theta^{\text{old}} - \alpha \nabla_\theta J(\theta)$$

\begin{itemize}
    \item $\theta$：模型的所有参数（对于 Word2Vec，就是所有的词向量）
    \item $\alpha$：学习率（step size），通常是一个很小的数（如 $10^{-3}$ 到 $10^{-5}$）
    \item $\nabla_\theta J(\theta)$：损失函数关于参数的梯度
\end{itemize}

\begin{figure}[H]
\centering
\includegraphics[width=0.85\textwidth]{slides-images/slide-004.jpg}
\caption{梯度下降更新公式：矩阵形式和逐参数形式\protect\footnotemark}
\end{figure}
\footnotetext{Slides 第5页。}

\begin{warningbox}{学习率不能太大}
如果学习率 $\alpha$ 过大，参数更新的步幅过大，可能会``跳过''最小值，甚至导致损失函数发散。因此必须选择足够小的学习率，使参数在每一步都朝着正确的方向微调。
\end{warningbox}

\subsection{随机梯度下降（SGD）}

标准梯度下降需要在\textbf{整个训练语料}上计算损失函数和梯度，这在数据量很大时极其耗时。实际中，我们使用\textbf{随机梯度下降}（Stochastic Gradient Descent, SGD）：

\begin{figure}[H]
\centering
\includegraphics[width=0.95\textwidth]{slides-images/slide-005.jpg}
\caption{随机梯度下降：每次只在一个小的数据子集（mini-batch）上计算梯度，大幅加速训练\protect\footnotemark}
\end{figure}
\footnotetext{Slides 第6页。}

\begin{importantbox}{SGD 的核心思想}
每次只选取一个小批量（mini-batch，如 16 或 32 个样本）数据，在这个子集上估计梯度，并据此更新参数。虽然每步的梯度估计是\textbf{有噪声的}，但：
\begin{enumerate}
    \item 训练速度大幅提升（不必等待遍历全部数据）
    \item 噪声反而有正则化效果，帮助模型跳出局部最优，实际上往往能得到更好的结果
\end{enumerate}
\end{importantbox}

\begin{lstlisting}[caption={随机梯度下降伪代码},language=Python]
while True:
    window = sample_window(corpus)
    theta_grad = evaluate_gradient(J, window, theta)
    theta = theta - alpha * theta_grad
\end{lstlisting}

\subsection{本章小结}

梯度下降是所有神经网络训练的基础。在实践中，我们几乎总是使用随机梯度下降（SGD）或其变体（如 Adam），因为它不仅比全量梯度下降快得多，而且由于引入的噪声，往往能找到更好的解。参数初始化必须使用\textbf{随机小数}（而非全零），否则会因为对称性无法学习。

%% ========== Part 2: Word2Vec 回顾与深入 ==========

